\documentclass[10pt,twocolumn]{article}
\usepackage[a4paper,margin=18mm,columnsep=6mm]{geometry}
\usepackage{booktabs,graphicx,amsmath,amssymb,xcolor,hyperref,multirow}
\usepackage[numbers,sort&compress]{natbib}
\newcommand{\Fone}{F$_1$}
\title{What Moves Detection and What Moves the Boundary?\\Target Source, Encoder Block and Fine-Tuning in Self-Supervised\\Detection of Killer-Whale Calls under Limited Data}
\author{Iaroslav Sheipak$^{1}$ and co-authors \\[2pt] \small $^1$MIPT}
\date{Draft v0.3 --- 19 October 2026}
\begin{document}
\maketitle

\begin{abstract}
We study which ingredients of a HuBERT-style self-supervised pipeline move event detection and which move boundary precision, on 17 field recordings of killer whales with 613 expert-annotated calls of 11 types and a pretraining corpus of 45k unlabelled lab clips. One encoder, one masking scheme, one decoder and one scorer are held fixed while the target source (log-mel $k$-means, random quantiser, $k$-means on a bioacoustic teacher, continuous EMA targets), the transformer block (vanilla vs.\ a Gemma-3-style stack with gated convolution), the token stride and the adaptation regime (frozen probe vs.\ full fine-tuning) are varied, with three seeds per main cell and paired bootstrap over recordings. Three results hold. (1) The effect of the block depends on the target type: with frozen probes it changes teacher-unit targets by $+0.011$ and continuous targets by $+0.19$ class-aware event \Fone{} at a 0.2\,s collar; after fine-tuning both lines reach 0.65--0.66. (2) After fine-tuning, 15--18M students pretrained for 30 minutes reach class-aware \Fone{} comparable to the fully fine-tuned 94M teacher (0.660$\pm$0.024 and 0.650$\pm$0.014 vs.\ 0.635$\pm$0.005), while the teacher remains ahead on the screening criterion that matters to annotators, class-agnostic recall at a fixed false-alarm budget (0.742 vs.\ 0.709 at 100 FA/h), and on onset precision. (3) A decomposition of the remaining misses shows that strict event-level losses are dominated by abutting calls in dense bouts rather than by sensitivity: 20\% of the references are found but scored as misses because one predicted event covers several calls, for teacher and students alike. A one-logit onset head added at fine-tuning recovers part of this ($+0.04$ to $+0.07$ strict recall at matched false-alarm rate), but a published onset-based detector, Voxaboxen, re-scored under the same protocol from stored predictions, is stronger at matched false-alarm rate. We also report audited negative results: event-level units built from oracle boundaries, segment units, dual codebooks, summed objectives, a second HuBERT iteration and block-masked JEPA. All conclusions are about this corpus and protocol; we do not claim a new architecture or generality across domains.
\end{abstract}

\section{Introduction}
Self-supervised audio and time-series models are evaluated mostly on clip classification or forecasting, while the task behind bioacoustic monitoring is to find every call in hours of recording and to place its onset and offset. A biologist screening hydrophone data wants high recall at a tolerable number of false alarms per hour; a benchmark wants event-level \Fone{} with a time collar. These two views of the same detector can disagree, and part of this paper is about where they disagree.

We make three contributions on one well-annotated corpus. First, a controlled comparison inside a HuBERT-style pipeline~\citep{hsu2021hubert}: the target source, the transformer block, the token stride and the adaptation regime are varied one at a time with everything else fixed, each main cell is run with three seeds, and differences are given with paired bootstrap intervals over recordings. Second, a direct comparison of compact students against their 94M teacher under both the benchmark metric and the screening criterion. Third, a decomposition of the misses of the best models that identifies dense bouts of abutting calls, not sensitivity, as the dominant strict-recall loss, a minimal mechanism that recovers part of it, and a re-scoring of a published onset-based detector under the same protocol. We also report, with the reasons, which natural ideas did not work: replacing frame-level units by event-level units even with oracle boundaries, segment units from teacher change-points, dual codebooks, summed discrete and continuous objectives, a second HuBERT iteration, and block-masked JEPA.

\section{Related work}
HuBERT~\citep{hsu2021hubert} predicts $k$-means units of masked frames; BEST-RQ~\citep{chiu2022bestrq} replaces $k$-means by a random projection quantiser; data2vec~\citep{baevski2022data2vec} regresses the latents of an EMA teacher; SUPERB~\citep{yang2021superb} fixed the weighted-layer frozen probe we adopt. AVES~\citep{hagiwara2023aves} ports HuBERT to animal sound and is our teacher. For event-level evaluation we use the collar-based metric of sed\_eval~\citep{mesaros2016metrics}. MarineMammalBench~\citep{sheipak2026mmb} introduced these recordings as a benchmark, reported frozen AVES heads at 39.7 event-\Fone{} under nested cross-validation with 1\,s windows, and derived the grid--tolerance ceiling $\mathrm{recall}(\tau)\le\min(1,2\tau/g)$ for tokens of stride $g$. Onset-based detectors predict event instances rather than frames: Voxaboxen~\citep{mahon2025voxaboxen} predicts, on a self-supervised backbone, whether a call starts in each window and how long it lasts, mirrors this from offsets and fuses both box sets; YOHO~\citep{venkatesh2022yoho} regresses start and end per class; sound event bounding boxes~\citep{ebbers2024sebbs} derive events from change points of frame-wise scores. Earlier supervised pipelines on the same recordings (a windowed CNN, ConvNeXt sliding windows, a Voxaboxen model trained on synthetic recordings and a two-stage hybrid) are listed in Table~\ref{tab:ft}.

\section{Task, metrics and protocol}
\label{sec:metric}
\paragraph{Task and scorer.} A recording contains events $(s_i,e_i,c_i)$. The model emits per-token probabilities over 11 call types and a noise class, a decoder turns them into events, and a prediction matches a reference if $|\hat s-s|\le\tau$, $|\hat e-e|\le\max(\tau,0.5(e-s))$ and, for the class-aware variant, the class agrees. This is the sed\_eval event-based metric with onset collar $\tau$; our greedy onset-nearest matching agreed with sed\_eval's optimal matching to within 0.1 \Fone{} point on every prediction set we checked. We sweep $\tau\in\{0.05,0.1,0.2,0.3,0.5,1\}$\,s and report class-aware \Fone{} at $\tau$=0.2 as the headline, \Fone{} at $\tau$=1\,s (called \Fone{}@1s; it is not a plateau: \Fone{} keeps rising by 0.03--0.12 up to $\tau$=10\,s), class-agnostic \Fone{} at $\tau$=0.2, and the median and 90th percentile of the absolute onset error of matched events (class-aware matching at $\tau$=1). Because matched sets differ between models, onset errors are also compared on the common set of references matched by all models under comparison.
\paragraph{Screening criterion.} Class-agnostic recall at a false-alarm budget: for each test fold, the two decoder thresholds are chosen on the two development recordings as the pair with the highest recall whose development false-alarm rate (FAR) does not exceed $B\in\{50,100,300\}$ per hour, and test recall is reported together with the \emph{actual} test FAR. Thresholds transfer imperfectly from two recordings, so test FAR can exceed the budget; where we compare models at equal load we read recall off a test recall--FAR curve obtained by sweeping the detection threshold. We also report coverage recall, the fraction of references with at least 50\% of their duration covered by some prediction, as a proxy for what an annotator reviewing the detections would see. Matching for the screening criterion uses $\tau$=0.5\,s.
\paragraph{Splits, seeds, intervals.} Three recording-held-out folds (6/6/5 recordings; the split of~\citep{sheipak2026mmb}); within each training set the first recording of each of the two training folds is held out as the development set; test results are pooled over the three test folds. Main cells have three pretraining seeds and are reported as mean$\pm$sd. Paired differences carry 95\% bootstrap intervals over the 17 recordings (2000 resamples); the recording is the resampling unit and carries all its seeds jointly. $P(A>B)$ denotes the share of positive bootstrap differences, not a $p$-value.

\section{Experimental setup}
\paragraph{Field data.} 17 continuous recordings (74\,min, 48\,kHz, Avacha Gulf) with 613 expert events of 11 call types (K1 103, K7 120, K5 87, K4 76, K17 63, K10 47, K27 39, K12 38, K13 22, K14 11, K21 7). Annotations of consecutive calls frequently abut or overlap.
\paragraph{Pretraining corpus.} 45{,}429 isolated lab clips (11.15\,h, 12{,}228 noise), no field audio. Each 5\,s example is a deterministic synthetic mixture: a cross-faded noise background with one to four faded calls at random gains (38\% call time on average). Audio is resampled to 16\,kHz; the student input is an 80-bin log-mel at 10\,ms hop stacked in pairs to 20\,ms tokens.
\paragraph{Encoders.} \emph{Vanilla}: 8 pre-norm transformer layers, $d$=384, 6 heads, FFN 1536, convolutional positional embedding, 14.8M parameters. \emph{Modern}: RMSNorm, GeGLU FFN, RoPE, QK-norm, local (64-token) and global attention alternating 3:1, and a Conformer gated-convolution module per layer; 17.8M parameters. Span masking with start probability 0.08 and span 10 tokens ($\approx$55\% masked); 20k steps of batch 32$\times$5\,s, AdamW, peak lr $5\cdot10^{-4}$, 1.5k warm-up, cosine decay; 30 minutes per run on an RTX 5090 laptop GPU.
\paragraph{Targets.} \textit{km(mel)}: $k$-means ($K$=100) on normalised stacked log-mel; \textit{BEST-RQ}: fixed random projection and codebook ($K$=100); \textit{km(AVES)}: $k$-means on layer-9 features of frozen AVES-bio (HuBERT-base, 94M, 20\,ms frames); \textit{data2vec}: regression of the instance-normalised mean of the top four layers of an EMA teacher; \textit{HuBERT iteration 2}: $k$-means on layer 6 of the km(AVES) student; \textit{JEPA}: a two-layer predictor with learned mask tokens over a context encoder that sees zeros at masked positions, regressing instance-normalised last-layer EMA latents, with block masks (mean 45 tokens, 36\% masked) or HuBERT-style short spans. Targets at 20\,ms are aligned token by token; the AVES front-end yields 499 tokens per 10\,s, and we pad each chunk to the exact grid, which matters for every number involving the frozen teacher (an earlier version of this work had a 20\,ms drift per chunk).
\paragraph{Adaptation and decoding.} Frozen probe: softmax-weighted layer sum, LayerNorm, linear 12-way head, 3k AdamW steps on 8k-token batches. Full fine-tuning: 1.5k steps of batch 8$\times$10\,s, encoder lr $5\cdot10^{-5}$ ($2\cdot10^{-5}$ for AVES), head lr $10^{-3}$. Decoder: $p_{\mathrm{call}}=1-p_{\mathrm{noise}}$, median filter of 5 tokens, two thresholds (detect, edge) selected on the development recordings, minimum duration 0.15\,s; the class of an event is the argmax of the summed call probabilities inside it.

\section{Results}
\begin{figure}[t]\centering\includegraphics[width=\linewidth]{f1_vs_tau.png}
\caption{Event-\Fone{} against onset collar $\tau$ (log scale), seed-0 runs. Solid and dashed: frozen probes with modern and vanilla blocks; dash-dotted: full fine-tuning. The curves differ mainly by vertical scale; after fine-tuning they also shift left.}
\label{fig:curves}\end{figure}

\begin{table*}[t]\centering\small
\caption{Frozen linear probes, 20\,ms tokens. Class-aware event-\Fone{} at $\tau$=0.05, 0.2 and 1\,s, class-agnostic \Fone{} at $\tau$=0.2, onset-error percentiles of matched events (ms). $n$ = number of pretraining seeds; sd is given where $n$=3. The frozen AVES row is the corrected (grid-aligned) evaluation.}
\label{tab:probe}
\begin{tabular}{llcccccc}\toprule
targets & block & $n$ & \Fone{}$_{\tau=.05}$ & \Fone{}$_{\tau=.2}$ & \Fone{}@1s & any$_{\tau=.2}$ & onset p50 / p90\\\midrule
log-mel, no context & -- & 1 & 0.156 & 0.268 & 0.391 & 0.351 & 56 / 318\\
random init & vanilla & 1 & 0.107 & 0.187 & 0.314 & 0.348 & 78 / 387\\
AVES-bio frozen (94M) & -- & 1 & 0.348 & 0.497 & 0.598 & 0.619 & 34 / 274\\\midrule
km(mel) & vanilla & 1 & 0.134 & 0.261 & 0.361 & 0.460 & 62 / 325\\
BEST-RQ & vanilla & 1 & 0.179 & 0.310 & 0.418 & 0.473 & 67 / 360\\
BEST-RQ & modern & 3 & 0.208 & 0.348$\pm$.038 & 0.458$\pm$.054 & 0.471 & 56 / 338\\
data2vec & vanilla & 3 & 0.139 & 0.275$\pm$.009 & 0.379$\pm$.010 & 0.407 & 68 / 379\\
data2vec & modern & 3 & 0.292 & \textbf{0.466$\pm$.031} & \textbf{0.581$\pm$.024} & \textbf{0.573} & \textbf{45 / 254}\\
km(AVES) & vanilla & 3 & 0.271 & 0.437$\pm$.027 & 0.571$\pm$.016 & 0.541 & 50 / 314\\
km(AVES) & modern & 3 & 0.268 & 0.448$\pm$.003 & 0.564$\pm$.009 & 0.539 & 52 / 340\\
HuBERT iteration 2 & vanilla & 3 & 0.245 & 0.422$\pm$.007 & 0.555$\pm$.010 & 0.518 & 53 / 322\\
JEPA, block masks & modern & 1 & 0.125 & 0.276 & 0.389 & 0.417 & 76 / 333\\
JEPA, short spans & modern & 1 & 0.215 & 0.415 & 0.537 & 0.530 & 60 / 270\\
\bottomrule\end{tabular}\end{table*}

\begin{table*}[t]\centering\small
\caption{Full fine-tuning, 20\,ms tokens. Columns as in Table~\ref{tab:probe}; the last column is sed\_eval overall event-\Fone{} ($t$\_collar 0.2, 50\% length) of the seed-0 run, for comparison with prior pipelines on the same recordings (lower block; only the ConvNeXt row is explicitly cross-validated, the hybrid used a field-fine-tuned classifier without reported cross-validation).}
\label{tab:ft}
\begin{tabular}{llcccccccc}\toprule
init & block & $n$ & \Fone{}$_{\tau=.05}$ & \Fone{}$_{\tau=.2}$ & \Fone{}@1s & any$_{\tau=.2}$ & onset p50 / p90 & sed\_eval (seed 0)\\\midrule
random & vanilla & 1 & 0.159 & 0.311 & 0.460 & 0.458 & 77 / 336 & --\\
random & modern & 3 & 0.296 & 0.464$\pm$.010 & 0.571$\pm$.004 & 0.603 & 42 / 286 & 45.9\\
km(AVES) & vanilla & 3 & 0.437 & 0.627$\pm$.015 & 0.729$\pm$.023 & 0.690 & 34 / 197 & 61.9\\
km(AVES) & modern & 3 & 0.498 & \textbf{0.660$\pm$.024} & \textbf{0.754$\pm$.020} & 0.715 & 29 / 171 & \textbf{67.8}\\
HuBERT iteration 2 & vanilla & 3 & 0.411 & 0.633$\pm$.014 & 0.732$\pm$.010 & 0.677 & 38 / 199 & --\\
BEST-RQ & modern & 3 & 0.426 & 0.590$\pm$.009 & 0.678$\pm$.012 & 0.673 & 33 / 226 & --\\
data2vec & modern & 3 & 0.477 & 0.650$\pm$.014 & 0.752$\pm$.008 & 0.715 & 32 / 194 & 65.7\\
JEPA, short spans & modern & 1 & 0.406 & 0.612 & 0.711 & 0.686 & 37 / 228 & --\\
AVES-bio (94M) & -- & 2 & \textbf{0.536} & 0.635$\pm$.005 & 0.721$\pm$.012 & \textbf{0.739} & \textbf{24 / 161} & 63.2\\\midrule
\multicolumn{10}{@{}p{0.97\textwidth}@{}}{\footnotesize\textit{Prior pipelines on the same 17 recordings (sed\_eval overall \Fone{}):} windowed CNN (ANIMAL-SPOT) 47.8; Perch 2.0 + linear probe 60.9; ConvNeXt sliding window, 4-fold CV 64.7; Voxaboxen, synthetic-trained 65.5; Voxaboxen + ConvNeXt hybrid 74.7; AVES full FT, nested CV, 1\,s windows~\citep{sheipak2026mmb} 41.2.}\\
\bottomrule\end{tabular}\end{table*}

\subsection{Result 1: the block matters for continuous targets, little for teacher units}
\label{sec:r1}
With frozen probes (Table~\ref{tab:probe}) the modern block changes km(AVES) by $+0.011$ (0.437$\to$0.448, inside seed spread) and data2vec by $+0.191$ (0.275$\to$0.466; seed-0 paired bootstrap $+0.161$ [$+0.072,+0.259$]); three seeds per cell agree. Unit-prediction accuracy of km(AVES) rises with the block (0.46$\to$0.52) without a downstream gain, so discrete teacher targets already supply the structure the block otherwise has to provide, whereas continuous EMA targets need it. After fine-tuning (Table~\ref{tab:ft}) the picture changes: the block helps km(AVES) by $+0.033$ (0.627$\to$0.660; paired bootstrap of seed-0 runs $+0.059$ [$+0.018,+0.094$]) and both lines end within 0.01 of each other (0.660$\pm$0.024 vs 0.650$\pm$0.014); the data2vec/vanilla cell was not fine-tuned, so the fine-tuned $2\times2$ is incomplete. The source of the targets orders the frozen-probe \Fone{}@1s values (0.36 km(mel), 0.42 BEST-RQ, 0.57 km(AVES)) in the same order as the mutual information of the units with the call types (PNMI 0.149 for km(mel), 0.272 for km(AVES) on the synthetic corpus), while all self-supervised probes place onsets within 45--68\,ms median; on the common set of 142 references matched by all models the onset medians are 36\,ms (frozen AVES), 51 (log-mel), 53 (km(AVES)) and 47 (data2vec, modern), and only fine-tuning moves them to 25--28\,ms. Pretraining under a fair comparison is worth $+0.19$ \Fone{} after fine-tuning over a scratch model with the same modern block and budget (0.464$\to$0.650; $+0.187$ [$+0.126,+0.254$]).

\subsection{Result 2: compact students versus the teacher}
\label{sec:r2}
Frozen, the 94M teacher is better than every student probe once its chunk timestamps are aligned correctly (0.497 vs 0.437--0.466; vs km(AVES) $+0.071$ [$-0.014,+0.148$]), and it localises better (34\,ms vs 45--53). Fine-tuned, the students reach comparable class-aware \Fone{} (km(AVES)/modern 0.660$\pm$0.024, data2vec/modern 0.650$\pm$0.014, teacher 0.635$\pm$0.005; km(AVES)/modern $-$ teacher $+0.046$ [$-0.005,+0.095$], seed-0 pair) at one fifth to one sixth of the parameters, but the teacher stays ahead on class-agnostic detection (0.739 vs 0.715) and onset precision (24 vs 29--32\,ms). The screening criterion (Table~\ref{tab:recall}) agrees with the latter: at a development budget of 100 FA/h the teacher reaches 0.742$\pm$0.014 recall (test FAR 84) and the best student 0.709$\pm$0.011 (test FAR 80); the modern block gives no recall gain over the vanilla HuBERT line (0.702 vs 0.702); a scratch model cannot meet the budget. With two teacher seeds and sd of 0.01--0.02 these differences are observed, not established.

\begin{table*}[t]\centering\small
\caption{Screening criterion: class-agnostic recall (mean$\pm$sd over seeds) at development budgets of 50/100/300 FA/h with the actual test FAR in brackets; full fine-tuning.}
\label{tab:recall}
\begin{tabular}{lcccc}\toprule
model & $n$ & R@50 (FAR) & R@100 (FAR) & R@300 (FAR)\\\midrule
AVES-bio FT & 2 & .744$\pm$.016 (69) & .742$\pm$.014 (84) & .753$\pm$.008 (103)\\
data2vec, modern & 3 & .716$\pm$.022 (63) & .709$\pm$.011 (80) & .728$\pm$.005 (114)\\
km(AVES), modern & 3 & .698$\pm$.030 (61) & .702$\pm$.025 (64) & .715$\pm$.022 (108)\\
km(AVES), vanilla & 3 & .697$\pm$.002 (65) & .702$\pm$.010 (79) & .719$\pm$.011 (124)\\
random, modern & 3 & .648$\pm$.004 (116) & .647$\pm$.004 (119) & .651$\pm$.017 (204)\\
\bottomrule\end{tabular}\end{table*}

\subsection{Result 3: where the strict-recall losses are}
\label{sec:r3}
\paragraph{Decomposition.} At the development-100 operating point (seed 0; test FAR 99 for the teacher, 79 for the data2vec student) the teacher detects 461 of 613 references (75\%) and the student 433 (71\%). Of the rest, 84/85 (14\%) are overlapped by a predicted event that fails the collar or offset test, 34/37 (6\%) are absorbed into a prediction already matched to a neighbour, 7/4 (1\%) are lost by the median filter or minimum duration, 15/30 (2/5\%) have a track peak between 0.1 and the threshold, and 12/24 (2/4\%) have no evidence in the track at all. Counting any overlap, 94.5\% (teacher) and 90.5\% (student) of the references are touched by a prediction; with the stricter $\ge$50\% coverage, 0.908 and 0.839 (three-seed means). The strict losses sit in dense bouts: 81\% of the boundary failures and 100\% of the absorptions have a neighbouring reference within 0.3\,s (median gap negative, i.e.\ overlapping annotations), against 15\% of the detected events; about 40\% of these neighbours are of a different type and 60\% of the same type. K12 accounts for 12\% of the teacher's misses. The teacher--student gap is elsewhere: 27 vs 54 references below threshold or absent, i.e.\ sensitivity.

\paragraph{A one-logit onset head.} We add one auxiliary logit per token during fine-tuning that predicts whether a call onset lies within $\pm$1 token (BCE, positive weight 20) and split each two-threshold event at peaks of the onset posterior (height and minimum distance selected on the development recordings), the previous piece ending where the next begins and pieces shorter than 0.15\,s discarded. A second variant splits where the median-filtered argmax over call types changes; it cannot represent two simultaneous calls. Table~\ref{tab:bouts} reports both on the same fine-tuned models (three seeds each; strict recall is class-agnostic at $\tau$=0.5; entries are means over seeds, boundary and absorption counts are mean counts per seed).

\begin{table*}[t]\centering\footnotesize
\caption{Decoder variants at the development-100 FA/h operating point, three seeds per model. Strict recall (class-agnostic, $\tau$=0.5) as mean$\pm$sd, actual test FAR, recall read from the test recall--FAR curve at 100 and 300 FA/h (``--'' where the curve does not reach the point), coverage recall ($\ge$50\%), and mean counts per seed of boundary-failed / absorbed references.}
\label{tab:bouts}
\begin{tabular}{llcccccc}\toprule
model & decoder & strict R & test FAR & R@100 & R@300 & cov$\ge$50\% & bnd / abs\\\midrule
student, data2vec/modern & base & 0.715$\pm$.012 & 71 & 0.730 & 0.748 & 0.839 & 70 / 35\\
 & class split & 0.745$\pm$.013 & 128 & -- & 0.788 & 0.833 & 58 / 27\\
 & onset split & \textbf{0.755$\pm$.012} & 96 & \textbf{0.777} & 0.805 & \textbf{0.854} & 54 / 33\\
 & both & 0.772$\pm$.010 & 133 & -- & 0.823 & 0.833 & 41 / 27\\\midrule
teacher, AVES-bio & base & 0.754$\pm$.008 & 76 & 0.759 & 0.770 & 0.908 & 77 / 38\\
 & class split & 0.804$\pm$.010 & 139 & -- & 0.822 & 0.880 & 52 / 28\\
 & onset split & \textbf{0.821$\pm$.004} & 84 & \textbf{0.836} & 0.848 & \textbf{0.914} & 43 / 34\\
 & both & 0.830$\pm$.013 & 140 & -- & 0.850 & 0.879 & 40 / 24\\\midrule
Voxaboxen (stored predictions) & thr.\ 0.5 & 0.816 & 61 & -- & -- & 0.938 & 33 / 56\\
\bottomrule\end{tabular}\end{table*}

\paragraph{Effect, source and limits.} Paired bootstrap over recordings: onset split vs the base decoder of the same model gives $+0.039$ [$+0.012,+0.079$] strict recall for the student and $+0.071$ [$+0.028,+0.111$] for the teacher (seeds 0--1); against plain fine-tuning without the auxiliary head, $+0.045$ [$-0.003,+0.107$] and $+0.078$ [$+0.025,+0.130$]; the auxiliary loss alone changes nothing ($+0.006$ [$-0.022,+0.035$], $+0.007$ [$-0.012,+0.023$]). Coverage recall moves by only $+0.015$ [$+0.007,+0.025$] and $+0.010$ [$+0.004,+0.016$]: the main effect is the separation of series that were already found, the coverage gain is small, and nothing here shows that previously invisible calls are recovered. Boundary failures fall by 23\% for the student (70$\to$54) and 44\% for the teacher (77$\to$43), absorptions are unchanged (35$\to$33, 38$\to$34): onset peaks resolve abutting calls, not true overlaps. Onset splitting adds false alarms (student 261$\to$350 and teacher 198$\to$221 summed over seeds), hence the FAR-matched columns; the class split adds $\approx$60 FA/h and lowers coverage, so its gain is not at matched FAR. By the pre-registered reading the effect is intermediate (between $+0.03$ and $+0.08$).

\paragraph{Voxaboxen under the same protocol.} We re-scored the stored per-recording predictions of a Voxaboxen model (BEATs backbone) trained only on synthetic recordings assembled from the lab clips, i.e.\ without field labels, fusing its forward and backward boxes by keeping backward boxes that do not overlap a forward box at IoU 0.5 (its own graph matching gives 65.5 sed\_eval \Fone{} against 62.6 with this fusion, so our re-scoring is conservative); the stored boxes are pre-filtered at detection probability 0.5, so no lower operating points are available. At its default threshold it reaches 0.816 strict recall at a test FAR of 61/h with coverage 0.938 (last row of Table~\ref{tab:bouts}), which our best configuration, the fine-tuned teacher with onset split, reaches only at 84 FA/h, while at $\approx$60 FA/h our fine-tuned models sit near 0.75. Its remaining misses are dominated by absorptions (56) rather than boundary failures (33), so in the configurations tested here true overlaps are not resolved by either approach. We draw two conclusions: the one-logit split is not a competitive detector and is not proposed as one; and the decomposition is the result, since the strict-recall losses of frame-wise detectors on this corpus, for a 94M teacher and 18M students alike, are dominated by abutting calls, for which an instance-level (onset-based) formulation is the appropriate tool.

\section{Ablations and negative results}
\paragraph{Token stride.} With BEST-RQ targets and 20k steps, \Fone{}$_{\tau=.2}$ is 0.347/0.310/0.225 at 10/20/40\,ms (onset p50 58/67/86\,ms). The 40\,ms loss is the grid ceiling; the 10\,ms gain disappears when compute is matched (10\,ms at 10k steps: 0.303, \Fone{}@1s 0.409 vs 0.418), and a 10\,ms data2vec model does not improve after fine-tuning either (0.616 vs 0.647, one seed, with the local-attention window halved in seconds).
\paragraph{Event-level units with oracle boundaries.} On the synthetic corpus the boundaries of inserted calls are known. Units averaged over the true event segments and units averaged over random segments with the same length multiset were trained with the same vanilla encoder and compute (two seeds each). Frozen probes: frame units 0.437$\pm$0.027, oracle 0.361$\pm$0.009, random 0.303$\pm$0.033; fine-tuned: 0.627$\pm$0.015, 0.542$\pm$0.033, 0.447$\pm$0.023; onset medians after fine-tuning 34/44/56\,ms. Alignment matters (oracle$-$random $+0.137$ [$+0.079,+0.206$], seed-0 pair after fine-tuning), but one unit per event loses more than alignment returns: the seed-mean gap frame$-$oracle is 0.086 after fine-tuning and the seed-0 bootstrap $+0.055$ [$-0.004,+0.105$] includes zero. This is a negative result for this construction of segment targets, not for event-level units in general.
\paragraph{Segment units from teacher change-points, dual codebooks, summed objectives.} Merging AVES frames by cosine change-points (mean 3.9 tokens) with a segment-start head: 0.412 (vanilla) and 0.403 with the modern block (best onset median of the frozen probes, 46\,ms, but \Fone{}@1s 0.527); segment starts fall near event boundaries with probability 0.39 against 0.25 inside, and spectral-flux change points do no better (0.44 vs 0.24). Predicting km(mel) and km(AVES) units jointly: 0.410. Cross-entropy on AVES units plus data2vec regression on the modern block: 0.459$\pm$0.039 over three seeds, equal to either alone with twice the spread.
\paragraph{Second HuBERT iteration, JEPA, random quantiser.} Re-clustering layer 6 of the km(AVES) student (twice the pretraining compute) gives 0.422$\pm$0.007 frozen and 0.633$\pm$0.014 fine-tuned against 0.437$\pm$0.027 and 0.627$\pm$0.015. The JEPA variant with block masks is below a scratch model after fine-tuning (0.392 vs 0.464); with HuBERT-style short spans and the same predictor it reaches 0.415 frozen and 0.612 fine-tuned (one seed), so most of the loss is the masking granularity, and the formulation differs from our data2vec in three ways at once (predictor, zeroed context, last-layer target), so no single factor is identified. BEST-RQ on the modern block reaches 0.348$\pm$0.038 and 0.590$\pm$0.009.

\section{Unit analysis}
On the synthetic corpus, purity/PNMI of units against call types are 0.639/0.149 for km(mel), 0.685/0.272 for km(AVES) and 0.694/0.294 for segment units. On the field recordings with grid-aligned features, km(AVES) units reach PNMI 0.44; they change at 52\% of tokens, somewhat more often at event boundaries (0.574 within $\pm$1 token of a boundary) than inside events (0.515). The probe weights of km(AVES) concentrate on layers 6--8. The onset-error distributions of matched events are nearly identical across self-supervised frozen probes and are moved by the modern block for continuous targets (p90 379$\to$254\,ms) and by fine-tuning (p50 to 24--34\,ms).

\section{A telemetry sanity run}
To check that the pipeline runs unchanged on multichannel input we applied it to 63 BeamNG laps (2.07\,h, 15 channels at 30\,Hz) with ten overlapping event classes derived from simulator flags and controls, under an inductive protocol (per-fold pretraining and normalisation on non-test laps, scratch baseline with the same modern block). Frozen probes are weak for every initialisation (0.10--0.13 at $\tau$=0.2). After fine-tuning (one seed), pretraining gives $+0.087$ over all ten classes (0.647 vs 0.560) but only $+0.026$ on the classes that are not a deterministic function of one input channel (ABS, TCS, ESC, impacts, airtime: 0.560 vs 0.534), inside the expected noise. We report this as a sanity run, not as a second domain result.

\section{Limitations}
One site and species; 74 labelled minutes; a 16\,kHz front-end that cuts the top 2\,kHz of the calls; thresholds chosen on two development recordings, which makes operating points transfer imperfectly; two teacher seeds; the fine-tuned data2vec/vanilla cell is missing; the Voxaboxen comparison uses a synthetic-trained model, stored predictions and an approximate fusion, and no field-fine-tuned Voxaboxen was run; the telemetry labels are partly definitional. Numbers are not comparable to~\citep{sheipak2026mmb} except through the sed\_eval column.

\section{Conclusion}
On this corpus and protocol: the transformer block is decisive for continuous self-supervised targets and secondary for teacher-unit targets under frozen probes, and the two lines converge after fine-tuning; compact students match the fine-tuned teacher on class-aware event \Fone{} but not on screening recall or onset precision; and the dominant strict-recall loss of frame-wise detectors is the separation of abutting calls in dense bouts, which a one-logit onset head partly recovers and which an onset-based instance detector handles better. Event-level units as a replacement for frame units, in the constructions we tested, do not.

\section*{Reproducibility}
Every run is stored with hardware fingerprint, code snapshot, data checksums, training curve and telemetry in an SQLite experiment database; hypotheses and verdicts are versioned alongside; per-recording predictions of all reported models are saved and the paired bootstraps are reproducible from them. Code, the synthetic-mixture recipe and checkpoints will be released.

\bibliographystyle{plainnat}
\begin{thebibliography}{12}
\bibitem[Hsu et al.(2021)]{hsu2021hubert} W.-N. Hsu et al. HuBERT: Self-supervised speech representation learning by masked prediction of hidden units. \emph{IEEE/ACM TASLP}, 2021.
\bibitem[Chiu et al.(2022)]{chiu2022bestrq} C.-C. Chiu et al. Self-supervised learning with random-projection quantizer for speech recognition. \emph{ICML}, 2022.
\bibitem[Baevski et al.(2022)]{baevski2022data2vec} A. Baevski et al. data2vec: A general framework for self-supervised learning in speech, vision and language. \emph{ICML}, 2022.
\bibitem[Hagiwara(2023)]{hagiwara2023aves} M. Hagiwara. AVES: Animal vocalization encoder based on self-supervision. \emph{ICASSP}, 2023.
\bibitem[Yang et al.(2021)]{yang2021superb} S.-W. Yang et al. SUPERB: Speech processing universal performance benchmark. \emph{Interspeech}, 2021.
\bibitem[Mesaros et al.(2016)]{mesaros2016metrics} A. Mesaros, T. Heittola, T. Virtanen. Metrics for polyphonic sound event detection. \emph{Applied Sciences}, 2016.
\bibitem[Mahon et al.(2025)]{mahon2025voxaboxen} L. Mahon, B. Hoffman et al. Voxaboxen: onset-based detection of overlapping bioacoustic events. \emph{DCASE Workshop}, 2025.
\bibitem[Venkatesh et al.(2022)]{venkatesh2022yoho} S. Venkatesh, D. Moffat, E. R. Miranda. You Only Hear Once: a YOLO-like algorithm for audio segmentation and sound event detection. \emph{Applied Sciences}, 2022.
\bibitem[Ebbers et al.(2024)]{ebbers2024sebbs} J. Ebbers, R. Haeb-Umbach, R. Serizel. Sound event bounding boxes. \emph{Interspeech}, 2024.
\bibitem[Sheipak et al.(2026)]{sheipak2026mmb} I. Sheipak et al. MarineMammalBench: Task, temporal resolution, and adaptation in marine audio transfer. \emph{under review}, 2026.
\end{thebibliography}
\end{document}
